
\documentclass{article}
\usepackage{colortbl}
\usepackage{makecell}
\usepackage{multirow}
\usepackage{supertabular}

\begin{document}

\newcounter{utterance}

\centering \large Interaction Transcript for game `cladder', experiment `full\_v1.5\_default', episode 9125 with selfrefine{-}Qwen3.5{-}9B.
\vspace{24pt}

{ \footnotesize  \setcounter{utterance}{1}
\setlength{\tabcolsep}{0pt}
\begin{supertabular}{c@{$\;$}|p{.15\linewidth}@{}p{.15\linewidth}p{.15\linewidth}p{.15\linewidth}p{.15\linewidth}p{.15\linewidth}}
    \# & \multicolumn{2}{c}{Player} && \multicolumn{2}{c}{Game Master} \\
    \hline

    \theutterance \stepcounter{utterance}  
    & & & \multicolumn{4}{p{0.6\linewidth}}{
        \cellcolor[rgb]{0.9,0.9,0.9}{
            \makecell[{{p{\linewidth}}}]{
                \texttt{\tiny{[P1$\langle$GM]}}
                \texttt{You are an expert in causal inference. The following question is not a typical commonsense query, but rather a meticulously designed question created by a professor specializing in causal inference, intended to assess the students' mastery of the course content.} \\
\texttt{Imagine a self{-}contained, hypothetical world with only the following conditions, and without any unmentioned factors or causal relationships: Kraz has a direct effect on pexu and rukz. Pexu has a direct effect on hwax. Hwax has a direct effect on rukz. Kraz is unobserved. For those who are not pexu and are not hwax, the probability of rukz is 19\%. For those who are not pexu and are hwax, the probability of rukz is 29\%. For those who are pexu and are not hwax, the probability of rukz is 20\%. For those who are pexu and are hwax, the probability of rukz is 31\%. For those who are not pexu, the probability of hwax is 16\%. For those who are pexu, the probability of hwax is 90\%. For those who are pexu, would it be less likely to see rukz if the individual was not pexu?} \\
\texttt{Start your answer with 'Yes' or 'No', followed by additional reasoning or evidence to support your explanation.} \\
            }
        }
    }
     \\ \\

    \theutterance \stepcounter{utterance}  
    & \multicolumn{4}{p{0.6\linewidth}}{
        \cellcolor[rgb]{0.561,0.737,0.561}{
            \makecell[{{p{\linewidth}}}]{
                \texttt{\tiny{[P1$\rangle$GM]}}
                \texttt{No} \\
\\ 
\texttt{**Reasoning:**} \\
\\ 
\texttt{To answer this question, we must compare the probability of observing **rukz** for an individual who is **pexu** (the observed state) against the probability of observing **rukz** for an individual who is **not pexu** (the counterfactual state).} \\
\\ 
\texttt{1.  **Calculate \$P(Rukz | Pexu=1)\$:**} \\
\texttt{    We are given the conditional probabilities for **rukz** based on the joint states of **pexu** and **hwax**:} \\
\texttt{    *   \$P(Rukz | Pexu=1, Hwax=1) = 0.31\$} \\
\texttt{    *   \$P(Rukz | Pexu=1, Hwax=0) = 0.20\$} \\
\\ 
\texttt{    We are also given the probability of **hwax** given **pexu**:} \\
\texttt{    *   \$P(Hwax=1 | Pexu=1) = 0.90\$} \\
\texttt{    *   Therefore, \$P(Hwax=0 | Pexu=1) = 1 {-} 0.90 = 0.10\$} \\
\\ 
\texttt{    Using the law of total probability:} \\
\texttt{    \$P(Rukz | Pexu=1) = (0.31 \textbackslash{}times 0.90) + (0.20 \textbackslash{}times 0.10)\$} \\
\texttt{    \$P(Rukz | Pexu=1) = 0.279 + 0.020 = 0.299\$ (or 29.9\%)} \\
\\ 
\texttt{2.  **Calculate \$P(Rukz | Pexu=0)\$:**} \\
\texttt{    We are given the conditional probabilities for **rukz** for those who are **not pexu**:} \\
\texttt{    *   \$P(Rukz | Pexu=0, Hwax=1) = 0.29\$} \\
\texttt{    *   \$P(Rukz | Pexu=0, Hwax=0) = 0.19\$} \\
\\ 
\texttt{    We are given the probability of **hwax** given **not pexu**:} \\
\texttt{    *   \$P(Hwax=1 | Pexu=0) = 0.16\$} \\
\texttt{    *   Therefore, \$P(Hwax=0 | Pexu=0) = 1 {-} 0.16 = 0.84\$} \\
\\ 
\texttt{    Using the law of total probability:} \\
\texttt{    \$P(Rukz | Pexu=0) = (0.29 \textbackslash{}times 0.16) + (0.19 \textbackslash{}times 0.84)\$} \\
\texttt{    \$P(Rukz | Pexu=0) = 0.0} \\
            }
        }
    }
    & & \\ \\

    \theutterance \stepcounter{utterance}  
    & & & \multicolumn{2}{p{0.3\linewidth}}{
        \cellcolor[rgb]{0.9,0.9,0.9}{
            \makecell[{{p{\linewidth}}}]{
                \texttt{\tiny{[GM$|$GM]}}
                \texttt{No} \\
            }
        }
    }
    & & \\ \\

    \theutterance \stepcounter{utterance}  
    & & & \multicolumn{2}{p{0.3\linewidth}}{
        \cellcolor[rgb]{0.9,0.9,0.9}{
            \makecell[{{p{\linewidth}}}]{
                \texttt{\tiny{[GM$|$GM]}}
                \texttt{yes} \\
            }
        }
    }
    & & \\ \\

    \theutterance \stepcounter{utterance}  
    & & & \multicolumn{2}{p{0.3\linewidth}}{
        \cellcolor[rgb]{0.9,0.9,0.9}{
            \makecell[{{p{\linewidth}}}]{
                \texttt{\tiny{[GM$|$GM]}}
                \texttt{game\_result = LOSE} \\
            }
        }
    }
    & & \\ \\

\end{supertabular}
}

\end{document}
